% Written by paper/cpc/scripts/corpus_crosscheck.py from paper/cpc/data/corpus_crosscheck.json.
\begin{table}[tbp]
\centering
\small
\caption{readcon-core and the independent reader over the conformance corpus and the CON 3.1 corpus.
Each entry of the CON 3.1 corpus manifest names acceptance with a frame count, or rejection with an error class and a line, and the basis of that outcome in the released CON 3.1 text (the released specification, 125 clauses): explicit when clauses decide the outcome and the class, class-open when clauses require the rejection and name no class, optional when a clause permits the outcome without requiring it, and open when no clause decides the file.
The corpus has 95 explicit, 21 class-open, 4 optional and 23 open entries.
The conformance manifest names acceptance, or rejection with an error class, and its entries count as explicit.
A reader matches an entry when it returns every value that the entry names.
These are the verdict, the frame count, the error class, the line of the error where the reader reports one, and the periodicity and the column-4 mask of the first frame where the entry names them.
The reader also warns of a time-unit override exactly on the entries that name one.
For \texttt{read\_con}, the rewrite of each file that its entry marks keeps the version, minor and units of every frame of major 3.
It writes a frame of major 1 or 2 as major 3, with each unit that the frame declares and the same fields, as the Python, Julia and Fortran writers do.
That rewrite keeps the minor of a frame that declares one or whose velocities carry a time label, which pins it to minor 0; a frame of version 3.0 with neither comes back at minor 1, and no marked file contains one.
A reader that differs from an explicit entry disagrees with the specification, and one that differs from any other entry departs from the reading that entry records.
readcon-core \readconCargoVersion{} puts no line number in its Python exceptions, so its lines are not compared; it warns of a declared time unit that it does not apply only when it converts a frame to ASE (\texttt{ConFrame.to\_ase}), and the check converts every frame that has velocities.
\texttt{read\_con} and \texttt{iter\_con} are two paths through readcon-core; the independent reader is pure Python written from the released text, and it follows the reading the manifest records where that text leaves a case open.
Field-identical counts the files that both \texttt{read\_con} and the independent reader accept with equal fields: coordinates, velocities and cell equal in every bit, and the same masks, identifiers, symbols, masses, version, minor, declared and applied time unit and periodicity.}
\label{tab:corpus-crosscheck}
\resizebox{\textwidth}{!}{%
\begin{tabular}{@{}lrrrrrrrrrrr@{}}
\toprule
 & & \multicolumn{3}{c}{\texttt{read\_con}} & \multicolumn{3}{c}{\texttt{iter\_con}} & \multicolumn{3}{c}{Independent} & \\
\cmidrule(lr){3-5}\cmidrule(lr){6-8}\cmidrule(lr){9-11}
File group & Files & Match & Disagree & Not followed & Match & Disagree & Not followed & Match & Disagree & Not followed & Field-identical \\
\midrule
Conformance, valid & 14 & 14 & 0 & 0 & 14 & 0 & 0 & 14 & 0 & 0 & 14/14 \\
Conformance, invalid & 8 & 8 & 0 & 0 & 8 & 0 & 0 & 8 & 0 & 0 & -- \\
CON 3.1 corpus, accept & 62 & 62 & 0 & 0 & 62 & 0 & 0 & 62 & 0 & 0 & 62/62 \\
CON 3.1 corpus, reject & 81 & 81 & 0 & 0 & 81 & 0 & 0 & 81 & 0 & 0 & -- \\
\midrule
All & 165 & 165 & 0 & 0 & 165 & 0 & 0 & 165 & 0 & 0 & 76/76 \\
\bottomrule
\end{tabular}}
\end{table}
